\begin{table}[H]
\rowcolors{1}{white}{white}
\centering
\tiny
\renewcommand{\arraystretch}{0.8}
\label{tab:unterkategorien}

\begin{tabular}{p{1cm} p{8.8cm} p{2.5cm}}
\toprule
\textbf{Code} & \textbf{induktive Unterkategorie} & \textbf{Belegstellen} \\
\midrule

\multicolumn{3}{l}{\textbf{K1 Qualitätsverständnis und ideale Metadaten}} \\
\midrule
K1.1 & Mindestvalidität und Best Practice & I1\_002--I1\_004 \\ \hline
K1.2 & Qualität statt Quantität & I3\_002--I3\_003 \\ \hline
K1.3 & Metadaten als Inhalts- und Nutzungsvorschau & I3\_005 \\ \hline
K1.4 & semantischer Graph als Idealform & I2\_003--I2\_006 \\
\midrule

\multicolumn{3}{l}{\textbf{K2 Vollständigkeit, Reichhaltigkeit und Feldrelevanz}} \\
\midrule
K2.1 & relative Vollständigkeit & I2\_026 \\ \hline
K2.2 & empfohlene Felder über Pflichtfelder hinaus & I1\_007--I1\_010; I3\_013 \\ \hline
K2.3 & fachlich vorhandene Informationen übernehmen & I2\_004--I2\_005 \\ \hline
K2.4 & relevante Felder statt maximale Feldanzahl & I3\_002--I3\_003 \\
\midrule

\multicolumn{3}{l}{\textbf{K3 Formale und syntaktische Qualität}} \\
\midrule
K3.1 & DCAT-AP.de-Mindestvalidität & I1\_002 \\ \hline
K3.2 & SHACL-Konformität als Basisdimension & I1\_027--I1\_029; I3\_015 \\ \hline
K3.3 & URI- und Vokabularkorrektheit & I3\_003; I3\_018 \\ \hline
K3.4 & Trennung von fehlendem Feld und syntaktischem Fehler & I3\_017 \\ \hline
K3.5 & XML/RDF-Validität und Datentypprüfung & I3\_018 \\
\midrule

\multicolumn{3}{l}{\textbf{K4 Semantische und inhaltliche Qualität}} \\
\midrule
K4.1 & nichtssagende oder zu kurze Titel & I2\_008 \\ \hline
K4.2 & Titel-Beschreibung-Redundanz & I2\_008 \\ \hline
K4.3 & regionale Begriffe und Fachsprache & I2\_009--I2\_010; I1\_043 \\ \hline
K4.4 & verständliche und erklärende Beschreibung & I2\_027; I3\_005 \\ \hline
K4.5 & inhaltliche Widersprüche zwischen Metadatenfeldern & I2\_011; I2\_047 \\ \hline
K4.6 & Keywords, Themenfeld und Kategorie als semantische Prüffelder & I1\_041; I3\_028 \\
\midrule

\multicolumn{3}{l}{\textbf{K5 Technische Zugänglichkeit, Aktualität und Nutzbarkeit}} \\
\midrule
K5.1 & Dead Links und Linkrot & I1\_015; I2\_011 \\ \hline
K5.2 & funktionierende Distributions-, Download- und Servicelinks & I1\_015; I2\_028 \\ \hline
K5.3 & Aktualität und Harvesting-Zeitpunkt & I1\_019; I1\_024--I1\_025 \\ \hline
K5.4 & Maschinenlesbarkeit und nicht-proprietäre Formate & I3\_021 \\
\midrule

\multicolumn{3}{l}{\textbf{K6 Interoperabilität und Linked-Data-Qualität}} \\
\midrule
K6.1 & URI statt Literal & I1\_020--I1\_022 \\ \hline
K6.2 & externe Referenzsysteme, Wikidata und GND & I2\_005; I1\_022 \\ \hline
K6.3 & kontrollierte Vokabulare und Thesauri & I2\_005; I2\_055; I3\_018 \\ \hline
K6.4 & Querverbindungen zwischen vergleichbaren Datensätzen & I2\_006; I2\_056 \\ \hline
K6.5 & Organisationsgraph und Vermeidung von Blank Nodes & I2\_029 \\
\midrule

\multicolumn{3}{l}{\textbf{K7 Kontextsensitivität und Fit for Purpose}} \\
\midrule
K7.1 & domänenspezifische Feldrelevanz & I1\_006; I1\_009--I1\_010 \\ \hline
K7.2 & Geodaten und Normierungsgrad & I1\_031--I1\_033; I3\_022 \\ \hline
K7.3 & föderale Vergleichbarkeit & I1\_040 \\ \hline
K7.4 & Stakeholder-spezifische Qualitätsfragen & I1\_047--I1\_049 \\ \hline
K7.5 & unterschiedliche Perspektiven von Bereitstellenden und Nutzenden & I1\_019 \\
\midrule

\multicolumn{3}{l}{\textbf{K8 Rohdatenbezug und Metadaten-Daten-Kongruenz}} \\
\midrule
K8.1 & Formatkongruenz zwischen Metadaten und Ressource & I2\_045; I3\_020 \\ \hline
K8.2 & Serviceverhalten und tatsächliche Nutzbarkeit & I2\_046 \\ \hline
K8.3 & Dateninhalt gegen Titel und Beschreibung prüfen & I2\_047 \\ \hline
K8.4 & Bounding Box und Flächendeckung aus Daten ableiten & I2\_013; I2\_059 \\ \hline
K8.5 & fachliche Grenzen der Rohdatenbewertung & I3\_022 \\
\midrule

\multicolumn{3}{l}{\textbf{K9 Scoring, Gewichtung, Reporting und Verbesserungshinweise}} \\
\midrule
K9.1 & wissenschaftlich begründete Metriken & I2\_021--I2\_024 \\ \hline
K9.2 & SHACL als Teildimension statt Gesamtscore & I1\_028--I1\_029 \\ \hline
K9.3 & Teil-Scores zur Problemlokalisierung & I2\_039--I2\_040 \\ \hline
K9.4 & Ranking als Verbesserungsanreiz & I2\_034--I2\_037 \\ \hline
K9.5 & verständliche Fehlermeldungen und Handlungshinweise & I2\_040--I2\_041; I3\_017 \\
\midrule

\multicolumn{3}{l}{\textbf{K10 KI-gestützte Analyseverfahren und Automatisierungsgrenzen}} \\
\midrule
K10.1 & Metadaten-Linter & I1\_039 \\ \hline
K10.2 & Freitextprüfung und Beschreibungsgüte & I1\_035; I1\_043--I1\_044 \\ \hline
K10.3 & Keyword-Mapping und Assisted Learning & I1\_041; I2\_055--I2\_056 \\ \hline
K10.4 & SPARQL- und Analyseunterstützung & I1\_047--I1\_049 \\ \hline
K10.5 & semantischer Graph als bessere KI-Grundlage & I2\_052--I2\_056 \\ \hline
K10.6 & Rohdatenanalyse durch KI & I2\_058--I2\_061; I3\_024 \\ \hline
K10.7 & Datenhoheit bei automatisierter Verbesserung & I2\_032; I3\_026 \\ \hline
K10.8 & Halluzination, stochastische Ausreißer und fachliche Fehlinterpretation & I1\_037; I1\_046; I2\_052 \\ \hline
K10.9 & Ressourcenaufwand und Betriebslast & I1\_046; I2\_062--I2\_063 \\

\bottomrule
\end{tabular}
\end{table}